%!TEX root = ../main.tex
\chapter{Sensor-Agnostic Perception Pipeline}
\label{ch:pipeline}

\section{Design Principles and System Overview}
\label{sec:pipeline-overview}

The pipeline described in this chapter performs the perception, state
estimation and environment representation functions on which the rest of this
work rests. The requirement that guided its design is platform independence:
the same system must operate on ground and aerial robots equipped with
different sensors, geometries and kinematics, without the code being rewritten.
On every platform it solves, in real time, the two problems that precede any
decision about motion: estimating the pose of the robot, and representing the
space that surrounds it.

The requirement is met through a uniform design principle: no node of the
system holds any information about the device that feeds it. A node knows the
type of data it receives and the processing it must apply to it, while the
identity of the producer is declared in the configuration. There are no
conditional branches on the model of a sensor; there is a configuration file
declaring which sensors are present, where they are mounted, on which channels
they publish and with which parameters, and the launch files build the system
accordingly, instantiating the nodes, connecting the channels and computing the
geometric transforms. It follows that the composition logic of the system
resides in the launch files rather than in the nodes, that introducing a sensor
of an already supported type requires no new code, and that the configuration is
generated by a dedicated tool rather than written by hand.

The architecture is organised in seven functional layers
(Table~\ref{tab:layers}), each with a single responsibility and communicating
with the others exclusively through the publish--subscribe channels provided by
the middleware. The separation allows one component to be replaced without
modifying the remaining ones. The same separation also provides the structure of
diagnosis, since the boundaries between layers are observable channels. The cost
of such an organisation is that every datum traverses a larger number of nodes,
with the resulting increase in latency and computational load.

\begin{table}[H]
  \centering
  \caption{The seven layers of the pipeline and their responsibilities.}
  \label{tab:layers}
  \begin{tabular}{@{}llp{7.8cm}@{}}
    \toprule
    & \textbf{Layer} & \textbf{Responsibility} \\
    \midrule
    L1 & Acquisition   & bring the datum into the system in a standardised form \\
    L2 & Preprocessing & validate and filter the datum, and declare its reliability \\
    L3 & TF            & publish the geometry of the robot \\
    L4 & Fusion        & combine and analyse the data so as to obtain one odometry and one observation of the surrounding world \\
    L5 & Map           & build the global map and the local map \\
    L6 & Planning      & decide the motion \\
    L7 & Actuation     & the checks that put the resulting command into effect \\
    \bottomrule
  \end{tabular}
\end{table}

In the present work the pipeline serves three distinct functions. It is the
producer of the observation the network consumes, defined in
Section~\ref{sec:observation-space}; it is the instrument with which the
training corpus described in Chapter~\ref{ch:training} was recorded; and it is
the environment in which the learnt policy is executed and evaluated. Layer L6
is the only one this work replaces with a learnt policy: it receives the
representation of the surrounding space, the pose and the goal, and publishes a
velocity command on the same channel on which a human operator would publish it.

Consistently with that role, the sections that follow describe the pipeline in
terms of its operating logic --- what each layer guarantees to the next, and for
what reason --- while detail is reserved for the observation space, which
constitutes the input contract of the network and conditions every choice made
in Chapter~\ref{ch:lnn}. Figure~\ref{fig:pipeline-chain} summarises the chain
they describe.

\begin{figure}[htbp]
  \centering
  \linespread{1}\selectfont
  \begin{tikzpicture}[x=1cm, y=1cm,
                      lyr/.append style={minimum width=17mm},
                      ext/.append style={minimum width=20mm}]
    \node[lyr] (l3) at (5.0,1.5) {L3\\TF};

    \node[ext] (win) at (0,0)   {simulator\\or hardware};
    \node[lyr] (l1)  at (2.4,0) {L1\\acquisition};
    \node[lyr] (l2)  at (4.4,0) {L2\\preproc.};
    \node[lyr] (l4)  at (6.4,0) {L4\\fusion};
    \node[lyr] (l5)  at (8.4,0) {L5\\map};
    \node[lyr] (l6)  at (10.4,0){L6\\policy};

    \draw[flow] (win) -- (l1);
    \draw[flow] (l1) -- (l2);
    \draw[flow] (l2) -- (l4);
    \draw[flow] (l4) -- (l5);
    \draw[flow] (l5) -- (l6);

    \draw[thin flow] (l3) -- (l4);
    \draw[thin flow] (l3) -| (l5);

    \node[lyr] (l7)  at (8.4,-2)  {L7\\actuation};
    \node[ext] (wout) at (0,-2)   {simulator\\or hardware};
    \draw[flow] (l6.south) |- (l7.east);
    \draw[flow] (l7.west) -- node[lbl,above]{\texttt{cmd\_vel}} (wout.east);

    \draw[dashed, black!45] (1.2,1.1) -- (1.2,-2.6);
    \node[note, anchor=north] at (1.2,-2.65) {the only boundary};
  \end{tikzpicture}
  \caption{The chain of guarantees of Section~\ref{sec:pipeline-architecture}.
           Each layer assures the next a property that allows it to disregard
           everything upstream: L4 receives a filtered cloud and does not know
           which device produced it, and L6 receives occupancy and does not know
           which engine built it. L3 publishes the common geometry rather than
           sitting on the path, and L7 makes the command observable without
           being interposed on it.}
  \label{fig:pipeline-chain}
\end{figure}

\section{Pipeline Architecture}
\label{sec:pipeline-architecture}

The chain may be read as a succession of guarantees: each layer assures the next
one a property that allows it to disregard everything that precedes it. It is
this chain of guarantees, rather than the sequence of nodes, that makes the
system independent of the platform.

The first layer guarantees \emph{uniformity of form}. Devices of the same type
do not necessarily expose the same streams: a depth camera may provide the point
cloud already computed or only the depth image, and a camera may provide
rectified or raw images. Since the subsequent layers are designed for a single
case, the first layer reconstructs whatever the device does not provide, so that
downstream every type of datum always appears in the same form. A distinct case
is kept separate from this mechanism, namely one in which a capability does not
exist at all on the platform under consideration: the system keeps explicit the
difference between a stream that can be derived and one that no node is able to
produce, so that the absence of a capability is not interpreted downstream as a
fault.

The second layer guarantees \emph{declared reliability}. Every datum is checked
for plausibility and continuity, is reduced to what is useful downstream, and
above all is accompanied by the judgement the system forms about it. That
judgement does not remain a diagnostic annotation: it determines the decision
uniquely --- publication, publication with increased uncertainty, or
rejection --- and is conveyed downstream in the covariance matrix that
accompanies the measurement. Without an explicit declaration of reliability the
subsequent layers would have no criterion by which to distinguish a trustworthy
measurement from a compromised one. The layer produces no estimate and does not
cross-check sensors: each node observes a single device, and any processing that
requires a comparison between sources belongs to the fusion layer.

The third layer, TF, guarantees a \emph{common geometry}. Since every device measures
in its own reference frame, comparing and fusing heterogeneous measurements
requires those frames to be related to that of the vehicle. The relations
between the body and the sensors are static and derive from the same mounting
declaration used to build the model of the robot in simulation. To them are
added the two dynamic relations prescribed by the conventional model adopted in
ROS~2: the one between the odometry origin and the body, continuous but subject
to drift, and the one between the world frame and the odometry origin, corrected
but discontinuous. The distinction has a direct operational consequence, since
short-term motion control relies on the continuous frame, whereas localisation
in the environment relies on the corrected one.

The fourth and fifth layers guarantee the \emph{unified products} --- one pose
and one representation of the environment --- and are treated respectively in
Section~\ref{sec:pipeline-fusion} and Section~\ref{sec:observation-space}, since
both contribute to defining what the network observes.

The seventh layer, finally, guarantees the \emph{observability of the command}.
It is not interposed on the path that carries the command to the actuators, but
republishes it time-stamped and at a constant rate, even in the absence of new
commands. It thereby becomes possible to determine which command was in force at
a given instant, a property on which the recording of the training corpus
depends (Section~\ref{sec:observation-recording}).

Two conventions, lastly, cut across all layers. The naming of the channels
declares the point of the chain at which the datum stands and the node that
produced it, and is constructed by a single shared module, since a divergence in
naming would not raise any detectable error, but the silent absence of
communication between two components. Uncertainty, as observed above,
constitutes the language in which sources are weighed against one another: from
that convention follows a property that conditions the design of the fusion
layer, namely that a source declaring an untruthful uncertainty degrades the
estimate more than its absence would.

\section{Multi-Sensor Data Fusion}
\label{sec:pipeline-fusion}

The fourth layer is the first in which measurements from different sensors are
compared. From that comparison derive the two unified products on which the rest
of the system rests: the pose, that is a single estimate of the state of the
vehicle, and a single point cloud expressed in the body frame, obtained by
temporally aligning the measurements of sources that differ in field of view,
range, density and rate.

The estimates produced by the individual sources are collected by a recursive
Bayesian filter, which maintains the state of the vehicle and corrects it at
every incoming measurement, weighing it by its own uncertainty. Since the
weighing rests on the declared covariance, the correctness of that declaration
directly conditions the quality of the estimate: many odometry implementations
leave the covariance field at zero, not because the estimate is exact but
because they do not compute its uncertainty, and the filter interprets the zero
value as certainty, anchoring itself to it and disregarding the remaining
sources. The layer therefore comprises a family of nodes whose task is to
attribute to the estimates the uncertainty they do not declare, modulating it
wherever a measurable quantity exists on which to base the attribution --- as in
the case of a geometrically degenerate environment, which makes one direction of
motion unobservable, or of wheel slippage, which invalidates the odometry
derived from it --- and declaring it constant elsewhere, since in the absence of
any evidence a declared constant is preferable to an arbitrary modulation.

It remains to be established, when several sources observe the same quantities,
which of them determines each component of the state. A criterion founded on
accuracy would require the sources to be ordered from the most precise to the
least precise, and that ordering to be recalibrated for every combination of
platform, terrain and environment, which would contradict the principle of
platform independence stated in Section~\ref{sec:pipeline-overview}. The system
therefore adopts a criterion founded on the \emph{type} of measurement, which is
a stable property of the device and not of the environment in which it operates:
direct velocity measurements are all admitted, since they do not compete with
one another and weighing them is the filter's task; external references not
subject to drift acquire ownership of the dimension they observe, and no
estimate obtained by integration may displace them; integrated estimates cover
the remaining dimensions, and those not elected contribute their own increments
rather than imposing their own pose. The configuration of the filter is
generated from that policy on the basis of the declared set of sensors, and
verified against a set of invariants that make structurally impossible the error
conditions the filter would otherwise report at run time.

\section{Observation Space}
\label{sec:observation-space}

This section defines the input of the network. It is the point at which the
pipeline ceases to be a perception system and becomes the contract on which
Chapter~\ref{ch:lnn} rests: every quantity listed here is a quantity the network
observes, and every representational choice made here is a choice the training
inherits.

\subsection{From the local map to the ego-centric crop}
\label{sec:observation-map}

The fifth layer builds two representations of the environment, distinct in
purpose and in requirements. The global map maintains the long-term consistency
of the space traversed and is used to correct the drift of the odometry through
the recognition of places already visited. The local map represents the space
immediately surrounding the vehicle as a grid of occupied voxels, and only the
latter reaches the network. The substantial difference lies in memory: the
global map must retain, whereas the local map must forget, since an obstacle
that has moved must disappear quickly and a cumulative representation would
populate the space with obstacles that no longer exist. It is a property
required by the use the network makes of it: the learnt policy reacts to what
surrounds the vehicle at the current instant, and a residual obstacle would be
indistinguishable from a real one.

The last node of the layer prepares the input of the network
(Figure~\ref{fig:observation}): the occupied voxels are brought into the body
frame of the robot and cropped within a sphere of fixed radius, bounded above
and below in height. Two choices deserve to be motivated. Expressing them in the
body frame keeps the coordinates bounded, preventing them from growing with the
distance travelled and making the observation invariant with respect to the
absolute position of the vehicle in the environment. The spherical shape, in
place of the rectangular volume in which the map is built, answers a requirement
of directional uniformity: in a rectangular volume the distance from the centre
depends on the direction, so that the network would receive a field of
observation varying with orientation, that is a systematic distortion that would
be learnt along with the behaviour. The crop is three-dimensional --- the radius
is measured in space, and the horizontal section therefore narrows towards the
top --- while the bounds in height remain independent of the radius, since the
space usable by a platform is limited in elevation.

\begin{figure}[htbp]
  \centering
  \linespread{1}\selectfont
  \begin{tikzpicture}[x=1cm, y=1cm]
    \draw[black!25, step=2.4mm] (-1.32,-1.32) grid (1.32,1.32);
    \draw[semithick] (-1.32,-1.32) rectangle (1.32,1.32);
    \draw[very thick] (0,0) circle (1.15);
    \fill (0,0) circle (1.1pt);
    \node[lbl, anchor=north] at (0,-1.5) {local map, quantised\\and cropped};

    \node[blk, anchor=west, text width=47mm] (coords) at (2.4,1.25)
      {occupied voxel indices\\[2pt]\footnotesize integer $(i,j,k)$, occupancy only};
    \node[blk, anchor=west, text width=47mm] (dense) at (2.4,-1.25)
      {dense observation vector\\[2pt]\footnotesize goal 7 $\mid$ odom 13 $\mid$ state 21};

    \draw[flow] (1.45,0.7) -- (coords.west);
    \draw[flow] (1.45,-0.7) -- (dense.west);

    \node[blk, minimum height=18mm, anchor=west] (net) at (8.6,0) {network};
    \draw[flow] (coords.east) -| ([xshift=-4mm]net.west) -- (net.west);
    \draw[flow] (dense.east)  -| ([xshift=-4mm]net.west) -- (net.west);

    \node[blk, anchor=west] (cmd) at (10.7,0) {\texttt{cmd\_vel}\\\footnotesize 6};
    \draw[flow] (net) -- (cmd);
  \end{tikzpicture}
  \caption{What the network observes. The local map is cropped into a sphere
           expressed in the body frame and quantised on a grid of the same pitch
           it was built with, reaching the network as the integer indices of the
           occupied cells --- occupancy alone, with no trace of the sensor that
           measured it. Beside it travels the dense vector of
           Table~\ref{tab:obs}, whose third block declares the body.}
  \label{fig:observation}
\end{figure}

\subsection{The representation of the cloud}
\label{sec:observation-cloud}

The cropped cloud is quantised on a grid aligned with the body axes, with the
same pitch used in the construction of the map, and delivered to the network as
a set of integer indices of the occupied voxels rather than as a set of
continuous coordinates. The representation is therefore purely binary: the
network observes which cells of the surrounding space are occupied, and neither
the density of the points that occupy them nor any photometric attribute. It
follows that the nature of the sensor that produced the measurement is not
observable by the network, which extends to the input of the policy the property
of platform independence that governs the whole pipeline.

The quantisation pitch and the crop radius are not redeclared at this level, but
read from the configuration of the nodes that actually applied them. The reason
is that quantising with a pitch different from the one with which the map was
built would produce no detectable error, but rather an input foreign to the
distribution on which the network was trained. The two parameters moreover vary
across platforms, and with them varies the extent of the grid: the corpus used
in this work comprises four of them, a circumstance the training accounts for by
stratifying the data split on the voxel pitch (Section~\ref{sec:sl-split}),
since a convolutional operator of fixed size covers a different physical extent
at a different resolution.

\subsection{The dense observation vector}
\label{sec:observation-vector}

Alongside the cloud the network receives a vector of \num{41} floating-point
values, organised in three blocks (Table~\ref{tab:obs}).

\begin{table}[H]
  \centering
  \caption{Composition of the dense observation vector.}
  \label{tab:obs}
  \begin{tabular}{@{}llp{7.2cm}@{}}
    \toprule
    \textbf{Block} & \textbf{Dim.} & \textbf{Content} \\
    \midrule
    Goal     & 7  & position and orientation of the assigned goal \\
    Odometry & 13 & position, orientation, linear velocity and angular velocity, from the fused pose \\
    State    & 21 & geometry and mass of the body (7), kinematic limits (12), voxel pitch (1), model identifier (1) \\
    \midrule
    \textbf{Total} & \textbf{41} & \\
    \bottomrule
  \end{tabular}
\end{table}

The first block is the goal, which the system receives from an operator or from
the mission configuration, expressed in the same frame as the pose: the network
therefore observes the goal and its own placement, and not the vector separating
them, whose construction is left to the network itself. The second block is the
pose produced by the fourth layer, with the linear and angular velocities that
accompany its estimate.

The third block is the element that makes a single network usable on a
heterogeneous fleet, and deserves to be examined in its components. The first
seven describe the body of the vehicle --- extent, centre of mass and mass ---
and thus place the dynamics of the robot within the observation. The following
twelve are the kinematic limits declared by the platform: minimum and maximum
velocity and acceleration on the linear, angular and vertical channels. They
constitute the scale with respect to which a command is to be interpreted, since
the same absolute velocity represents a moderate command on one platform and a
command at the limit on another; their role extends beyond the observation, in
that the loss function of supervised training uses them to make comparable the
errors committed on different robots (Section~\ref{sec:sl-loss}). A null limit,
moreover, declares that the corresponding axis does not exist on the platform,
as vertical velocity on a ground vehicle: the observation is therefore able to
represent the absence of a degree of freedom without the structure of the vector
changing.

The last two components declare the voxel pitch of the map --- which informs the
network of the resolution at which the concurrent cloud was built --- and the
identifier of the robot model. The latter is a categorical variable rather than
a physical quantity, and the treatment it receives inside the network is
discussed in Section~\ref{sec:aux-modules}.

The output is a six-component velocity command --- three linear and three
angular velocities --- expressed in the physical units of the platform and
published on the command channel of the vehicle, the same on which an operator
would publish. Nothing along this path saturates the command against the limits
declared in the third block: limiting, where necessary, is the task of the
on-board controller, and the network is evaluated on what it actually commands.

\subsection{Recording the corpus}
\label{sec:observation-recording}

The three components that consume the contract just defined --- the node that
records the corpus, the node that runs inference on board and the interface
towards the reinforcement learning environment --- share a single
implementation. The network is in fact trained on what the first produces and
executed inside the other two, and distinct implementations of the same
definition would diverge without raising any error, since the dimensions of the
vectors would continue to match and the only observable effect would be
erroneous behaviour.

Recording pairs in time the observation available to the vehicle with the action
that followed from it. The pairing is not symmetric: the pose is sought in the
neighbourhood of the instant of the observation, since it describes the state
the vehicle was in, whereas the command is sought \emph{afterwards}, since it
constitutes the reaction to it. The command considered is the one republished by
the seventh layer, whose re-emission at a constant rate makes it determinable
which command was in force at each instant, even when the operator has issued no
new ones. Every sample is moreover stamped with its own instant of acquisition:
the interval between consecutive samples is not constant, and constitutes in
turn an input of the network, for the reasons set out in
Section~\ref{sec:lnn-taxonomy}.

\section{Simulation--Hardware Parity}
\label{sec:pipeline-parity}

The system is designed so that simulation and real hardware are
interchangeable. The boundary between the two cases crosses the chain at a
single point, constituted by the boundary channels of the vehicle: on them
publishes the bridge towards the simulator when the robot is simulated, and the
driver of the device when it is real, while the command traverses the same
frontier in the opposite direction. Downstream of that point the nodes are the
same and hold no information about the origin of the data.

The two cases differ in the naming of the boundary channels alone: in simulation
it is established by the project, whereas with real hardware it is imposed by
the manufacturer of the device, in the absence of a common convention among
vendors. The difference is absorbed at a single point of the configuration, in
which the channel on which each stream actually arrives is declared. Moving to
real hardware therefore entails a change in three elements only --- the
component that publishes on the boundary channels, the source of time, and the
generated configuration --- while the seven layers, the frame tree, the tuning
parameters of the nodes and the fusion policy remain unchanged.

The property is essential to this work in two respects. The first concerns the
data: the training corpus is recorded through the same chain that will feed the
network during execution, so that the observation on which the policy is trained
and the one it receives on board are produced by the same code. The second
concerns geometry: since the model used in simulation and the static transforms
derive from the same mounting declaration, the geometry assumed by the system
and the one realised by the simulator coincide by construction, and the
discrepancy between simulation and reality does not include, at least, the
placement of the sensors.
